% MBBO205 Bulut Bilişim ve Yapay Zekâ -- Hafta 1 Ders Notu
% Derleme: xelatex MBBO205_Hafta01_Ders_Notu.tex (iki kez)
\documentclass[11pt,a4paper]{article}

\usepackage{fontspec}
\usepackage{polyglossia}
\setdefaultlanguage{turkish}
\setmainfont{TeX Gyre Pagella}
\setsansfont{TeX Gyre Heros}
\setmonofont{DejaVu Sans Mono}[Scale=0.85]

\usepackage[margin=2.3cm]{geometry}
\usepackage{amsmath,amssymb}
\usepackage{graphicx}
\usepackage{xcolor}
\usepackage{booktabs,tabularx,array}
\usepackage{enumitem}
\usepackage{fvextra}
\usepackage[most]{tcolorbox}
\usepackage{titlesec}
\usepackage{fancyhdr}
\usepackage[hidelinks]{hyperref}

\definecolor{lacivert}{HTML}{14213D}
\definecolor{turuncu}{HTML}{A94E17}
\definecolor{kum}{HTML}{EDEBE4}
\definecolor{kod}{HTML}{1B2A4A}
\definecolor{gri}{HTML}{4A5568}

\titleformat{\section}{\sffamily\Large\bfseries\color{lacivert}}{\thesection}{0.8em}{}
\titleformat{\subsection}{\sffamily\large\bfseries\color{lacivert}}{\thesubsection}{0.6em}{}
\setlength{\parindent}{0pt}
\setlength{\parskip}{0.55em}
\setlist{itemsep=0.2em,topsep=0.3em}
\renewcommand{\arraystretch}{1.25}
\newcolumntype{Y}{>{\raggedright\arraybackslash}X}

\pagestyle{fancy}
\fancyhf{}
\fancyhead[L]{\small\sffamily MBBO205 · Hafta 1}
\fancyhead[R]{\small\sffamily Dr. Öğr. Üyesi Hasan Oğuz}
\fancyfoot[C]{\small\sffamily\thepage}
\renewcommand{\headrulewidth}{0.4pt}

\newtcolorbox{kutu}[1][]{colback=kum,colframe=kum,arc=2mm,boxrule=0pt,left=3mm,right=3mm,top=2mm,bottom=2mm,#1}
\newtcolorbox{uyari}[1][]{colback=white,colframe=turuncu,arc=2mm,boxrule=0.8pt,left=3mm,right=3mm,top=2mm,bottom=2mm,title={#1},coltitle=white,colbacktitle=turuncu,fonttitle=\sffamily\bfseries}
\DefineVerbatimEnvironment{kodblok}{Verbatim}{breaklines,fontsize=\small,frame=single,rulecolor=\color{gri},framesep=2mm}

\begin{document}

%-------------------------------------------------------------- Başlık
\begin{tcolorbox}[colback=lacivert,colframe=lacivert,arc=2mm,left=5mm,right=5mm,top=4mm,bottom=4mm]
\color{white}\sffamily
{\small İSTANBUL OKAN ÜNİVERSİTESİ · MESLEK YÜKSEKOKULU · 2026–2027 GÜZ}\\[3mm]
{\LARGE\bfseries MBBO205 Bulut Bilişim ve Yapay Zekâ}\\[1mm]
{\large Hafta 1 Ders Notu: Derse giriş, yapay zekâya genel bakış ve Linux'a ilk adım}\\[3mm]
{\small Dr. Öğr. Üyesi Hasan Oğuz}
\end{tcolorbox}

\begin{kutu}
\textbf{Bu haftanın hedefleri.} Hafta sonunda öğrenci (i) yapay zekâ, makine öğrenmesi, derin öğrenme ve veri bilimi kavramlarını ayırt eder; (ii) alanın tarihsel gelişimini ve bugünkü sıçramanın nedenlerini açıklar; (iii) YZ sistemlerinin çalıştığı donanım ve bulut modellerini tanır; (iv) bir Ubuntu ortamında terminali açar, dizin yapısında gezinir, temel komutları kullanır ve ortam raporunu üretir. \\[1mm]
\textbf{İlgili kazanımlar:} DÖK1, DÖK2 \quad \textbf{Süre:} 1 saat teori + 1 saat uygulama
\end{kutu}

\tableofcontents

%==============================================================
\section{Ders bilgileri}

\subsection{İşleyiş}
Ders haftada bir ders saati teori ve bir ders saati Linux uygulamasından oluşur. Her hafta aynı döngü izlenir:
\begin{enumerate}
  \item \textbf{Teori:} kavramlar ve gerekli matematik.
  \item \textbf{Uygulama:} aynı konunun Linux terminalinde yapılması.
  \item \textbf{Teslim:} hafta içinde, 5. haftadan itibaren otomatik denetimli.
  \item \textbf{Bozuk sistem görevi:} kasten bozulmuş bir ortamı onarıp kök nedeni raporlama.
  \item \textbf{İngilizce okuma:} ilgili resmî dokümantasyon veya \texttt{man} sayfası.
\end{enumerate}

\subsection{Değerlendirme}
\begin{center}
\begin{tabularx}{0.92\textwidth}{@{}l c Y@{}}
\toprule
\textbf{Bileşen} & \textbf{Ağırlık} & \textbf{Kapsam} \\
\midrule
Ara sınav & \%40 & Uygulamalı sınav (Bash, SQL, veri keşfi; hafta 1–7) ve hafta 1–7 lab teslimleri \\
Final & \%60 & Kavramsal sorular ve bozuk sistem onarımı; dönem projesi ve hafta 8–14 lab teslimleri \\
\bottomrule
\end{tabularx}
\end{center}
Sınav notlarının iç dağılımı dönem başında Telegram grubunda duyurulur.

\subsection{İletişim ve ders grubu}
\begin{minipage}[c]{0.62\textwidth}
Duyurular, haftalık lab dosyaları ve soru-cevap için ders Telegram grubuna katılın:\\[2mm]
\texttt{[ders grubu bağlantısı]}\\[2mm]
E-posta gönderirken konu satırına \texttt{[MBBO205]} yazın; ad, soyad ve öğrenci numaranızı ekleyin. Görüşme saatleri: [gün, saat].
\end{minipage}\hfill
\begin{minipage}[c]{0.32\textwidth}
\end{minipage}

\subsection{Ders kuralları}
\begin{itemize}
  \item \textbf{Akademik dürüstlük:} birlikte düşünün, ayrı teslim edin.
  \item \textbf{Üretken YZ:} kullanımı serbesttir; kullanılan istemler ve doğrulama adımları \texttt{AI\_USAGE.md} dosyasında beyan edilir. Ara sınav ve final kapalı ağ ortamında yapılır.
  \item \textbf{Terminal güvenliği:} \texttt{sudo} ve \texttt{rm -r} komutlarını, ne yaptığını okumadan çalıştırmayın.
\end{itemize}

%==============================================================
\section{Yapay zekâ nedir?}

\subsection{Tanımlar}
\textbf{Yapay zekâ (YZ)}, normalde insan zekâsı gerektiren görevleri (algılama, dil, akıl yürütme, karar verme, deneyimden öğrenme) yerine getiren hesaplamalı sistemlerin bilimi ve mühendisliğidir. Mühendislik açısından en kullanışlı tanım \textbf{rasyonel ajan} tanımıdır: çevresini algılayan ve hedefine göre beklenen en iyi sonucu veren eylemi seçen sistem \cite{aima}. Bu tanım zekânın ``ne olduğunu'' değil, sistemin \emph{davranışını} ölçer.

Bugünkü sistemler \textbf{dar YZ}'dir: belirli görevlerde (görüntü tanıma, çeviri, oyun) çok güçlüdürler. İnsan düzeyinde genel yeteneğe sahip \textbf{genel YZ (AGI)} ise tanımı ve ulaşılabilirliği hâlâ tartışmalı bir hedeftir.

\subsection{Kavramlar arasındaki ilişki}
Kavramlar iç içe kümeler oluşturur:
\[
\text{Üretken YZ} \subset \text{Derin öğrenme} \subset \text{Makine öğrenmesi} \subset \text{Yapay zekâ}.
\]
\begin{itemize}
  \item \textbf{Yapay zekâ:} öğrenmeyen yöntemleri de kapsar; kural tabanlı uzman sistemler, arama (satranç motorları), planlama.
  \item \textbf{Makine öğrenmesi (ML):} davranışını açıkça programlanmak yerine veriden öğrenen modeller.
  \item \textbf{Derin öğrenme (DL):} çok katmanlı yapay sinir ağlarıyla yapılan makine öğrenmesi.
  \item \textbf{Üretken YZ:} yeni metin, görüntü, ses veya kod üreten derin modeller (büyük dil modelleri, LLM).
  \item \textbf{Veri bilimi:} bu kümelerin \emph{içinde değil}, onlarla \emph{kesişen} bir alandır; istatistik, programlama ve alan bilgisini birleştirerek veriden karar üretir.
\end{itemize}

\subsection{Programlama paradigması olarak makine öğrenmesi}
Geleneksel programlamada insan kuralları yazar; program kurallarla veriyi işleyip cevap üretir. Makine öğrenmesinde ise veri ve doğru cevaplar verilir; öğrenme algoritması kuralları (modeli) çıkarır:
\begin{align*}
\text{Geleneksel:}&\quad \text{Kurallar} + \text{Veri} \;\longrightarrow\; \text{Cevaplar}\\
\text{Makine öğrenmesi:}&\quad \text{Veri} + \text{Cevaplar} \;\longrightarrow\; \text{Kurallar (model)}
\end{align*}
Örnek: bir spam filtresi için binlerce kural yazmak yerine, ``spam'' ya da ``normal'' olarak etiketlenmiş binlerce e-posta verilir; model hangi sözcük örüntülerinin spam'e işaret ettiğini kendisi öğrenir.

\subsection{Öğrenme türleri}
\begin{center}
\begin{tabularx}{\textwidth}{@{}l Y Y@{}}
\toprule
\textbf{Tür} & \textbf{Ne öğrenir?} & \textbf{Örnek} \\
\midrule
Denetimli & Etiketli veriden girdi $\to$ çıktı eşlemesi & Spam tespiti, ev fiyatı tahmini \\
Denetimsiz & Etiketsiz veride yapı (kümeler, aykırılar) & Müşteri segmentasyonu, anomali tespiti \\
Pekiştirmeli & Ödül ve cezayla deneme-yanılma politikası & AlphaGo, robot kontrolü \\
Öz-denetimli & Etiketi verinin kendisinden üretir & LLM'lerin sonraki sözcüğü tahmin etmesi \\
\bottomrule
\end{tabularx}
\end{center}
Dönem boyunca ağırlıklı olarak denetimli öğrenmeyle çalışılacaktır (hafta 9–11).

%==============================================================
\section{Tarihçe}

\begin{center}
\begin{tabularx}{\textwidth}{@{}l Y@{}}
\toprule
\textbf{Yıl} & \textbf{Olay ve önemi} \\
\midrule
1943 & McCulloch ve Pitts, yapay nöronun ilk matematiksel modelini önerir \cite{mcculloch}. \\
1950 & Turing, ``Makineler düşünebilir mi?'' sorusunu taklit oyunuyla ele alır \cite{turing}. \\
1956 & Dartmouth çalıştayı; ``yapay zekâ'' terimi alanın adı olur \cite{dartmouth}. \\
1958 & Rosenblatt'ın algılayıcısı (perceptron): veriden öğrenen ilk yapay nöron \cite{rosenblatt}. \\
1966 & ELIZA: ilk sohbet programlarından biri. \\
1969 & Minsky ve Papert, tek katmanlı ağların sınırlarını gösterir (ör. XOR). \\
1974–1980 & Birinci ``YZ kışı'': beklentiler karşılanamaz, fonlar kesilir. \\
1980'ler & Uzman sistemler: elle yazılmış kural tabanları ticari olarak yaygınlaşır. \\
1986 & Geri yayılım (backpropagation) çok katmanlı ağları eğitilebilir kılar \cite{rumelhart}. \\
1987–1993 & İkinci YZ kışı: uzman sistem pazarı çöker. \\
1997 & Deep Blue, dünya satranç şampiyonu Kasparov'u yener. \\
2006 & Hinton ve ekibi derin ağları katman katman ön eğitimle eğitir. \\
2009 & ImageNet: milyonlarca etiketli görüntüden oluşan veri seti. \\
2012 & AlexNet, GPU ile eğitilen bir CNN ile ImageNet'te hatayı büyük ölçüde düşürür \cite{alexnet}. \\
2016 & AlphaGo, Lee Sedol'u Go'da 4–1 yener \cite{alphago}. \\
2017 & ``Attention Is All You Need'': Transformer mimarisi \cite{transformer}. \\
2020 & GPT-3 (175 milyar parametre) \cite{gpt3}; AlphaFold 2 protein yapısı tahmininde çığır açar \cite{alphafold}. \\
2022 & ChatGPT ile üretken YZ geniş kitlelere ulaşır. \\
2024 & Fizik Nobeli Hopfield ve Hinton'a, Kimya Nobeli Hassabis, Jumper ve Baker'a \cite{nobel}; AB YZ Yasası yürürlüğe girer \cite{aiact}. \\
2025+ & YZ ajanları, açık ağırlıklı ve yerel çalışan modeller. \\
\bottomrule
\end{tabularx}
\end{center}

\subsection{Algılayıcı: ilk öğrenen nöron}
Rosenblatt'ın algılayıcısı, girdilerin ağırlıklı toplamını bir eşikle karşılaştırır:
\begin{equation}
y = \begin{cases} 1, & \mathbf{w}^{\top}\mathbf{x} + b > 0,\\ 0, & \text{aksi halde,}\end{cases}
\qquad \mathbf{w}^{\top}\mathbf{x} = \sum_{j=1}^{n} w_j x_j .
\label{eq:perceptron}
\end{equation}

\begin{kutu}
\textbf{Çözümlü örnek: VE kapısı.} $n=2$, $\mathbf{w}=(1,1)$ ve $b=-1.5$ seçelim. Dört girdi için:
\begin{center}
\begin{tabular}{cc|c|c}
$x_1$ & $x_2$ & $w_1x_1+w_2x_2+b$ & $y$ \\ \hline
0 & 0 & $-1.5$ & 0 \\
0 & 1 & $-0.5$ & 0 \\
1 & 0 & $-0.5$ & 0 \\
1 & 1 & $\phantom{-}0.5$ & 1
\end{tabular}
\end{center}
Algılayıcı mantıksal VE işlevini gerçekler. XOR için ise $(0,1)$ ve $(1,0)$ noktalarını $(0,0)$ ve $(1,1)$ noktalarından ayıran tek bir doğru yoktur; bu nedenle tek bir algılayıcı XOR'u öğrenemez. Minsky ve Papert'in 1969 eleştirisinin özü budur. Çok katmanlı ağlar bu sınırı aşar.
\end{kutu}

\subsection{Neden şimdi?}
Algoritmaların önemli bölümü onlarca yıllıktır. 2012 sonrasındaki sıçramayı üç etkenin aynı anda olgunlaşması açıklar:
\begin{enumerate}
  \item \textbf{Veri:} internet, akıllı telefonlar ve sensörler; ImageNet gibi büyük etiketli veri setleri.
  \item \textbf{Hesaplama:} GPU'lar, sinir ağlarının temel işlemi olan matris çarpımını binlerce çekirdekte paralel yapar ve bulutta saatlik kiralanabilir.
  \item \textbf{Algoritma ve açık kaynak:} geri yayılım, ReLU, dropout, Transformer; PyTorch ve Hugging Face gibi araçların herkese açık olması.
\end{enumerate}
YZ kışlarının ortak nedeni tersidir: vaatler, o dönemin veri ve hesaplama olanaklarının önüne geçmiştir.

%==============================================================
\section{Alt disiplinler ve yaşam döngüsü}

\subsection{Alt disiplinler}
\begin{center}
\begin{tabularx}{\textwidth}{@{}l Y@{}}
\toprule
\textbf{Alan} & \textbf{Tipik görevler} \\
\midrule
Görüntü işleme (CV) & Sınıflandırma, nesne tespiti, bölütleme, tıbbi görüntü analizi \\
Doğal dil işleme (NLP) & Çeviri, duygu analizi, özetleme, soru-cevap, büyük dil modelleri \\
Konuşma & Konuşmadan metne (ASR), metinden konuşmaya (TTS) \\
Robotik & Algı (CV) + planlama + kontrol; otonom araçlar, endüstriyel robotlar \\
Öneri sistemleri & E-ticaret, video ve müzik platformlarında kişiselleştirme \\
\bottomrule
\end{tabularx}
\end{center}
Bu alanlar giderek ortak bir çekirdeği paylaşır: aynı derin öğrenme araçları ve çoğu zaman aynı Transformer mimarisi. \textbf{Çok kipli (multimodal)} modeller görüntü, metin ve sesi tek modelde işleyerek alanlar arasındaki sınırı belirsizleştirir. Robotik ise tanımı gereği çok disiplinlidir: bir robot kamerasıyla görür (CV), sözlü komutu anlar (NLP/konuşma) ve hareketini planlar (kontrol).

\subsection{Bir YZ projesinin yaşam döngüsü}
Aşağıdaki döngü dersin haritasıdır; parantez içindeki sayılar ilgili haftalardır.
\begin{enumerate}
  \item Problem tanımı (1)
  \item Veri toplama (3, 7)
  \item Temizleme ve ETL (7–8)
  \item Keşifsel veri analizi (6)
  \item Modelleme (9–12)
  \item Değerlendirme (10)
  \item Dağıtım (14)
  \item İzleme; yeni veri ve yeni sorunlarla başa dönüş (14)
\end{enumerate}
Bu bir çizgi değil döngüdür. Sektörde zamanın büyük bölümü 2. ve 3. adımlara harcanır.

%==============================================================
\section{Matematiksel öz}

Denetimli bir model, parametreleri $\theta$ olan bir fonksiyondur:
\begin{equation}
\hat{y} = f(\mathbf{x};\,\theta).
\end{equation}
Burada $\mathbf{x}\in\mathbb{R}^{n}$ girdi özellikleri, $y$ gerçek cevap, $\hat{y}$ modelin tahminidir. \textbf{Öğrenme}, $N$ örnekten oluşan eğitim verisi $\{(\mathbf{x}_i,y_i)\}_{i=1}^{N}$ üzerinde ortalama kaybı en aza indiren parametreleri bulmaktır:
\begin{equation}
\theta^{*} = \arg\min_{\theta}\; \frac{1}{N}\sum_{i=1}^{N} \mathcal{L}\big(f(\mathbf{x}_i;\theta),\, y_i\big).
\label{eq:erm}
\end{equation}
Kayıp fonksiyonu $\mathcal{L}$, tahminin ne kadar yanlış olduğunu ölçer. Örneğin sayısal tahmin (regresyon) için ortalama kare hata:
\begin{equation}
\mathrm{MSE}(\theta) = \frac{1}{N}\sum_{i=1}^{N}\big(y_i - \hat{y}_i\big)^{2}.
\end{equation}

\begin{kutu}
\textbf{Çözümlü örnek.} Üç evin gerçek fiyatları $y=(2,\,3,\,5)$ (milyon TL), bir modelin tahminleri $\hat{y}=(2.5,\,3,\,4)$ olsun. Hatalar $(-0.5,\,0,\,1)$, kareleri $(0.25,\,0,\,1)$ olur ve
\[
\mathrm{MSE} = \tfrac{1}{3}(0.25+0+1) \approx 0.417 .
\]
Eğitim, bu değeri küçültecek biçimde $\theta$'yı değiştirmektir.
\end{kutu}

Denklem~\eqref{eq:erm} dönemin matematiksel iskeletidir: lineer cebir $f$'nin nasıl hesaplandığını (hafta 5), olasılık ve istatistik verinin ve kaybın doğasını (hafta 6), optimizasyon ise $\arg\min$'in nasıl bulunduğunu (hafta 10–11) açıklar.

%==============================================================
\section{Altyapı: donanım ve bulut}

\subsection{Hangi iş hangi donanımda?}
\begin{center}
\begin{tabularx}{\textwidth}{@{}l Y Y@{}}
\toprule
\textbf{Birim} & \textbf{Güçlü olduğu} & \textbf{YZ'deki rolü} \\
\midrule
CPU & Az sayıda güçlü, genel amaçlı çekirdek & Veri hazırlama, küçük modeller, servis \\
GPU & Binlerce basit çekirdek; paralel matris çarpımı & Derin öğrenme eğitimi ve çıkarımı \\
TPU / NPU & Matris işlemine özel yonga (ASIC) & Büyük ölçekli eğitim; telefonda çıkarım \\
Uç cihaz & Düşük güç, yerel çalışma & Kamera, IoT, gizlilik gerektiren çıkarım \\
\bottomrule
\end{tabularx}
\end{center}
\textbf{Eğitim} pahalıdır ve seyrek yapılır; \textbf{çıkarım} (eğitilmiş modelin kullanılması) ucuzdur ama milyonlarca kez tekrarlanır. Donanım seçimi bu ayrıma göre yapılır.

\subsection{Bulut hizmet modelleri}
\begin{itemize}
  \item \textbf{IaaS (altyapı):} sanal makine, disk, ağ kiralanır; işletim sistemi ve üstü kullanıcıdadır. Bu derste hafta 4 ve 14'te kullanılacak.
  \item \textbf{PaaS (platform):} yönetilen veritabanı, model eğitim servisleri; kod ve veri kullanıcıda, sunucu sağlayıcıdadır.
  \item \textbf{SaaS (yazılım):} sohbet asistanları, e-posta, ofis uygulamaları; yalnızca kullanılır.
\end{itemize}
Bulut sunucularının büyük çoğunluğu Linux çalıştırır ve TOP500 listesindeki süper bilgisayarların tamamı Kasım 2017'den beri Linux kullanır \cite{top500}. IaaS seçen biri fiilen bir Linux sistem yöneticisidir; dersin Linux ağırlığının gerekçesi budur.

%==============================================================
\section{Güncel eğilimler ve riskler}

\subsection{Eğilimler}
\begin{enumerate}
  \item \textbf{Üretken ve çok kipli YZ:} metin, görüntü, ses ve kod üretimi.
  \item \textbf{YZ ajanları:} araç kullanan, web'de gezinen ve terminalde komut çalıştıran modeller.
  \item \textbf{Küçük ve yerel modeller:} açık ağırlıklı modeller dizüstü bilgisayarda çalıştırılabilir (hafta 13).
  \item \textbf{Uç YZ:} telefon ve kameralarda NPU ile yerel çıkarım.
  \item \textbf{Düzenleme:} AB YZ Yasası (Tüzük 2024/1689), KVKK, GDPR.
\end{enumerate}

\begin{uyari}{Öngörü (spekülatif)}
Ajanlar kod yazmayı kolaylaştırdıkça, bir sistemi kurma, denetleme ve sorun giderme becerisinin göreli değeri artabilir. Bu bir tahmindir; dersteki ``bozuk sistem'' görevleri bu varsayıma dayanır.
\end{uyari}

\subsection{Riskler}
\begin{itemize}
  \item \textbf{Önyargı (bias):} model, verideki dengesizliği öğrenir ve büyütebilir (hafta 10).
  \item \textbf{Gizlilik:} kişisel veriler KVKK ve GDPR kapsamındadır; anonimleştirme gerekir (hafta 7).
  \item \textbf{Halüsinasyon:} dil modelleri akıcı ama yanlış cevaplar üretebilir; çıktı doğrulanmalıdır (hafta 13).
  \item \textbf{Güvenlik:} istem enjeksiyonu, açık bırakılan portlar, kod deposuna sızan API anahtarları (hafta 4, 13).
\end{itemize}

%==============================================================
\section{Uygulama: Linux terminaline ilk adım}

\subsection{Kurulum seçenekleri}
Ders boyunca \textbf{Ubuntu 24.04 LTS} kullanılır. Önerilen en düşük donanım 8 GB RAM ve 40 GB boş disktir; GPU gerekmez.
\begin{itemize}
  \item \textbf{WSL2 (Windows):} Windows içinde gerçek bir Linux çekirdeği çalıştırır; en hızlı başlangıçtır \cite{wsl}.
  \item \textbf{VirtualBox:} macOS veya WSL'nin açılamadığı bilgisayarlar için sanal makine.
  \item \textbf{Lab bilgisayarı:} okulda hazır ortam. \textbf{Bulut VM:} hafta 4'ten itibaren.
\end{itemize}

\begin{kodblok}
# PowerShell (yönetici olarak), ardından bilgisayarı yeniden başlatın
wsl --install -d Ubuntu-24.04

# Ubuntu terminalinde: sürümü doğrulayın ve sistemi güncelleyin
lsb_release -a
sudo apt update && sudo apt upgrade -y

# WSL'de Windows dosyalarınız burada görünür:
ls /mnt/c/Users/
\end{kodblok}
Parola yazarken ekranda karakter görünmemesi normaldir.

\subsection{Bir komut satırının anatomisi}
\begin{kodblok}
ogrenci@okan-lab:~$ ls -la /etc
\end{kodblok}
\begin{center}
\begin{tabularx}{\textwidth}{@{}l Y@{}}
\toprule
\textbf{Parça} & \textbf{Anlamı} \\
\midrule
\texttt{ogrenci@okan-lab} & Kullanıcı adı @ makine adı \\
\texttt{\textasciitilde} & Bulunulan dizin; \texttt{\textasciitilde} ev dizinidir (\texttt{/home/ogrenci}) \\
\texttt{\$} & Normal kullanıcı; \texttt{\#} görülürse root (tam yetki) olarak çalışılıyordur \\
\texttt{ls} & Komut \\
\texttt{-la} & Seçenekler: \texttt{l} ayrıntılı liste, \texttt{a} gizli dosyalar dahil \\
\texttt{/etc} & Argüman: komutun üzerinde çalıştığı hedef \\
\bottomrule
\end{tabularx}
\end{center}
Linux büyük-küçük harfe duyarlıdır: \texttt{Veri.csv} ile \texttt{veri.csv} farklı dosyalardır. Boşluk içeren adlar tırnak içine alınmalıdır.

\subsection{Dizin yapısı (FHS)}
Windows'taki \texttt{C:}, \texttt{D:} sürücüleri yerine Linux'ta her şey tek bir kökten, \texttt{/} dizininden dallanır \cite{fhs}.
\begin{center}
\begin{tabularx}{\textwidth}{@{}l Y Y@{}}
\toprule
\textbf{Dizin} & \textbf{İçerik} & \textbf{YZ bağlamında örnek} \\
\midrule
\texttt{/home} & Kullanıcı dosyaları & Projeler, veri setleri \\
\texttt{/etc} & Sistem yapılandırması & \texttt{/etc/os-release}, SSH ayarları \\
\texttt{/usr} & Programlar ve kütüphaneler & \texttt{/usr/bin/python3} \\
\texttt{/var} & Değişken veri, günlükler & \texttt{/var/log}, veritabanı dosyaları \\
\texttt{/opt} & Üçüncü taraf yazılımlar & Bazı ticari araç kurulumları \\
\texttt{/tmp} & Geçici dosyalar & İndirmeler, ara çıktılar \\
\texttt{/proc} & Çekirdeğin sanal dosyaları & \texttt{/proc/cpuinfo}, \texttt{/proc/meminfo} \\
\texttt{/dev} & Aygıt dosyaları & Diskler, GPU aygıtları \\
\bottomrule
\end{tabularx}
\end{center}
``Linux'ta her şey bir dosyadır'' ilkesi gereği işlemci bilgisi bile \texttt{cat /proc/cpuinfo} ile okunabilir.

\subsection{Komut başvuru tablosu}
\begin{center}
\begin{tabularx}{\textwidth}{@{}l Y l@{}}
\toprule
\textbf{Komut} & \textbf{Ne yapar} & \textbf{Örnek} \\
\midrule
\texttt{pwd} & Bulunulan dizini yazar & \texttt{pwd} \\
\texttt{ls} & Dizin içeriğini listeler & \texttt{ls -la \textasciitilde} \\
\texttt{cd} & Dizin değiştirir & \texttt{cd ..}, \texttt{cd -} \\
\texttt{mkdir} & Dizin oluşturur & \texttt{mkdir -p a/b/c} \\
\texttt{touch} & Boş dosya oluşturur & \texttt{touch notlar.txt} \\
\texttt{cp} & Kopyalar & \texttt{cp -r veri/ yedek/} \\
\texttt{mv} & Taşır, yeniden adlandırır & \texttt{mv eski.txt yeni.txt} \\
\texttt{rm} & Siler (geri dönüşüm kutusu yok) & \texttt{rm -r dizin/} \\
\texttt{cat}, \texttt{less} & Dosyayı yazar, sayfa sayfa okur & \texttt{less env\_report.sh} \\
\texttt{head}, \texttt{tail} & İlk / son satırlar & \texttt{tail -n 5 log.txt} \\
\texttt{wc -l} & Satır sayar & \texttt{wc -l veri.csv} \\
\texttt{echo} & Metni yazar & \texttt{echo \$PATH} \\
\texttt{>}, \texttt{>>} & Çıktıyı dosyaya yazar / ekler & \texttt{echo a >> b.txt} \\
\texttt{|} & Çıktıyı sonraki komuta verir & \texttt{ls -S /etc | head} \\
\texttt{man}, \texttt{-{}-help} & Kılavuz, kısa yardım & \texttt{man ls} \\
\bottomrule
\end{tabularx}
\end{center}
\textbf{Kısayollar:} \texttt{Tab} tamamlama; \texttt{↑} ve \texttt{Ctrl+R} komut geçmişi; \texttt{Ctrl+C} çalışan komutu durdurur (kopyalamaz; terminalde kopyala-yapıştır \texttt{Ctrl+Shift+C/V}); \texttt{Ctrl+L} ekranı temizler.

\subsection{Sistem bilgisi komutları}
\begin{kodblok}
whoami               # hangi kullanıcıyım?
uname -r             # çekirdek sürümü
cat /etc/os-release  # dağıtım ve sürüm
lscpu | head -n 15   # işlemci modeli, çekirdekler
nproc                # kullanılabilir çekirdek sayısı
free -h              # bellek (RAM)
df -h /              # kök dizindeki disk alanı
nvidia-smi           # NVIDIA GPU varsa model ve bellek
\end{kodblok}
Çekirdek sayısı paralel veri işlemeyi, RAM belleğe sığacak veri seti ve model boyutunu, disk ise veri setleri ve model dosyalarını sınırlar. \texttt{nvidia-smi} ``command not found'' veriyorsa makinede NVIDIA sürücüsü veya GPU yoktur; bu bir hata değil, bir bilgidir.

%==============================================================
\section{Laboratuvar görevi}

\subsection{Ana görev: ortam raporu (yaklaşık 25 dakika)}
\begin{enumerate}
  \item Ders dizinini oluşturun ve içine girin.
  \item \texttt{env\_report.sh} dosyasını Telegram grubundan indirip bu dizine kopyalayın.
  \item Çalıştırma izni verin, çalıştırın ve çıktıyı inceleyin.
  \item Betiği açın; her satırın ne yaptığını kendi cümlelerinizle not edin.
\end{enumerate}
\begin{kodblok}
mkdir -p ~/mbbo205/hafta01
cd ~/mbbo205/hafta01
cp /mnt/c/Users/ADINIZ/Downloads/env_report.sh .   # WSL kullanıyorsanız
chmod +x env_report.sh
./env_report.sh env_report.txt
less env_report.sh
\end{kodblok}

Betiğin içeriği:
\begin{kodblok}
#!/usr/bin/env bash
# Hafta 1: ortam raporu -> env_report.txt
set -euo pipefail
out="${1:-env_report.txt}"
{
  echo "Tarih   : $(date -Iseconds)"
  echo "Kullanıcı: $(whoami)"
  echo "Kernel  : $(uname -r)"
  echo "Dağıtım : $(. /etc/os-release && echo "$PRETTY_NAME")"
  echo "CPU     : $(lscpu | awk -F: '/Model name/ {gsub(/^ +/,"",$2); print $2; exit}')"
  echo "Çekirdek: $(nproc)"
  echo "Bellek  : $(free -h | awk '/^Mem:/ {print $2}')"
  echo "Disk /  : $(df -h / | awk 'NR==2 {print $4 " boş / " $2}')"
  if command -v nvidia-smi >/dev/null; then
    echo "GPU     : $(nvidia-smi --query-gpu=name,memory.total --format=csv,noheader)"
  else
    echo "GPU     : yok"
  fi
  echo "Python  : $(python3 --version 2>&1 || echo yok)"
} | tee "$out"
\end{kodblok}

\begin{uyari}{Sık karşılaşılan hata: \texttt{bad interpreter}}
Betik Windows'ta bir metin düzenleyiciyle açılıp kaydedildiyse satır sonları CRLF olur ve Bash betiği çalıştıramaz. Düzeltme: \texttt{sed -i 's/\textbackslash r\$//' env\_report.sh}. Bu durum 3. haftanın bozuk sistem görevinde ayrıntılı ele alınacaktır.
\end{uyari}

\subsection{Bozuk sistem görevi: kayıp PATH (yaklaşık 15 dakika)}
Bir arkadaşınız terminalde şunu yazmış ve artık hiçbir komut çalışmıyor:
\begin{kodblok}
$ export PATH=/opt/yok
$ ls
ls: command not found
$ python3
python3: command not found
\end{kodblok}
\textbf{Görev:} sorunu aşağıdaki sırayla çözün ve \texttt{KOK\_NEDEN\_RAPORU.md} şablonunu doldurun.
\begin{enumerate}
  \item \textbf{Belirti:} hata mesajını aynen kaydedin.
  \item \textbf{Hipotezler:} en az iki olası neden yazın (ör. komut silinmiş olabilir; kabuk komutu yanlış yerde arıyor olabilir).
  \item \textbf{Testler:} \texttt{echo \$PATH} ve \texttt{/usr/bin/ls} komutlarını deneyin. Tam yolla çalışan komut ne anlatıyor?
  \item \textbf{Düzeltme:} PATH'i onarın.
  \item \textbf{Kök neden ve önleme:} 3--5 cümle.
\end{enumerate}
\textbf{İleri soru:} Aynı \texttt{export} satırı \texttt{\textasciitilde/.bashrc} dosyasının sonuna yazılmış olsaydı, yeni bir terminal açmak sorunu çözer miydi? Nasıl düzeltirdiniz?

\subsection{Teslim listesi}
\begin{itemize}
  \item \texttt{env\_report.txt}
  \item \texttt{KOK\_NEDEN\_RAPORU.md} (bozuk PATH görevi)
  \item \texttt{env\_report.sh} için satır satır açıklamalar
  \item \texttt{AI\_USAGE.md} (üretken YZ kullandıysanız)
\end{itemize}
Son teslim: [tarih, saat] \quad Teslim yeri: [LMS / e-posta]

%==============================================================
\section{Kendinizi sınayın}
\begin{enumerate}
  \item Bir satranç motoru ile bir spam filtresi arasındaki fark, ``YZ'' ve ``ML'' kavramları açısından nedir?
  \item Veri bilimi neden YZ kümesinin içinde değil de onunla kesişen bir alan olarak tanımlanır?
  \item Denklem~\eqref{eq:perceptron}'deki algılayıcı için VEYA kapısını gerçekleyen bir $(\mathbf{w},b)$ bulun.
  \item AlexNet'in 2012'deki başarısında GPU'nun rolü nedir? Aynı ağ 1990'da neden eğitilemezdi?
  \item $y=(1,\,4)$ ve $\hat{y}=(2,\,2)$ için MSE değerini hesaplayın.
  \item Bir üniversite, öğrenci sohbet asistanını IaaS, PaaS ve SaaS olarak nasıl kurabilir? Her birinde kim neyi yönetir?
  \item \texttt{echo "a" > x.txt} komutu iki kez çalıştırılırsa \texttt{x.txt} kaç satır içerir? \texttt{>>} ile kaç satır?
  \item \texttt{/usr/bin/ls} çalışıyor ama \texttt{ls} çalışmıyorsa sorun nerededir?
\end{enumerate}

\section*{Hafta 2 için hazırlık}
Ubuntu kurulumunu tamamlayın; \texttt{man hier} sayfasını okuyun; MIT ``The Missing Semester'' dersinin ``The Shell'' bölümünü izleyin \cite{missing}. Gelecek hafta: kullanıcılar, izinler, \texttt{apt}, \texttt{systemd} ve SSH.

%==============================================================
\begin{thebibliography}{99}\small
\bibitem{aima} S. Russell, P. Norvig, \emph{Artificial Intelligence: A Modern Approach}, 4. baskı, Pearson, 2020. \url{https://aima.cs.berkeley.edu/}
\bibitem{mcculloch} W. S. McCulloch, W. Pitts, ``A logical calculus of the ideas immanent in nervous activity,'' \emph{Bull. Math. Biophys.} 5, 115–133 (1943). \url{https://doi.org/10.1007/BF02478259}
\bibitem{turing} A. M. Turing, ``Computing machinery and intelligence,'' \emph{Mind} 59, 433–460 (1950). \url{https://doi.org/10.1093/mind/LIX.236.433}
\bibitem{dartmouth} J. McCarthy ve ark., ``A Proposal for the Dartmouth Summer Research Project on Artificial Intelligence'' (1955). \url{http://jmc.stanford.edu/articles/dartmouth/dartmouth.pdf}
\bibitem{rosenblatt} F. Rosenblatt, ``The perceptron: a probabilistic model for information storage and organization in the brain,'' \emph{Psychol. Rev.} 65, 386–408 (1958). \url{https://doi.org/10.1037/h0042519}
\bibitem{rumelhart} D. E. Rumelhart, G. E. Hinton, R. J. Williams, ``Learning representations by back-propagating errors,'' \emph{Nature} 323, 533–536 (1986). \url{https://doi.org/10.1038/323533a0}
\bibitem{alexnet} A. Krizhevsky, I. Sutskever, G. E. Hinton, ``ImageNet classification with deep convolutional neural networks,'' NeurIPS 2012. \url{https://papers.nips.cc/paper/4824-imagenet-classification-with-deep-convolutional-neural-networks}
\bibitem{alphago} D. Silver ve ark., ``Mastering the game of Go with deep neural networks and tree search,'' \emph{Nature} 529, 484–489 (2016). \url{https://doi.org/10.1038/nature16961}
\bibitem{transformer} A. Vaswani ve ark., ``Attention is all you need,'' NeurIPS 2017. \url{https://arxiv.org/abs/1706.03762}
\bibitem{gpt3} T. Brown ve ark., ``Language models are few-shot learners,'' NeurIPS 2020. \url{https://arxiv.org/abs/2005.14165}
\bibitem{alphafold} J. Jumper ve ark., ``Highly accurate protein structure prediction with AlphaFold,'' \emph{Nature} 596, 583–589 (2021). \url{https://doi.org/10.1038/s41586-021-03819-2}
\bibitem{nobel} Nobel Prize in Physics 2024. \url{https://www.nobelprize.org/prizes/physics/2024/summary/}
\bibitem{aiact} Regulation (EU) 2024/1689 (Artificial Intelligence Act). \url{https://eur-lex.europa.eu/eli/reg/2024/1689/oj}
\bibitem{top500} TOP500, Operating System istatistikleri. \url{https://www.top500.org/statistics/list/}
\bibitem{wsl} Microsoft, ``Install Linux on Windows with WSL.'' \url{https://learn.microsoft.com/en-us/windows/wsl/install}
\bibitem{fhs} Linux Foundation, \emph{Filesystem Hierarchy Standard 3.0}. \url{https://refspecs.linuxfoundation.org/FHS_3.0/fhs/index.html}
\bibitem{missing} MIT, \emph{The Missing Semester of Your CS Education}, ``Course overview + the shell.'' \url{https://missing.csail.mit.edu/2020/course-shell/}
\end{thebibliography}

\end{document}
